% Part: incompleteness
% Chapter: incompleteness-provability
% Section: lob-thm

\documentclass[../../../include/open-logic-section]{subfiles}

\begin{document}

\olfileid{inc}{inp}{lob}

\olsection{د لوب قضيه}

د يوې تيورۍ~$\Th{T}$ د ګوډل جمله د~$\lnot
\OProv[\Th{T}](y)$ يو ثابت ټکی دے؛ يعنې داسې !!a{sentence}~$!G$
چې
\[
\Th{T} \Proves \lnot \OProv[\Th{T}](\gn{!G}) \liff !G.
\]
دا جمله !!{derivable} نۀ ده۔ ځکه که~$\Th{T} \Proves !G$ وي، نو
(الف) د !!{derivability} د شرط~(۱) له مخې
$\Th{T} \Proves \OProv[\Th{T}](\gn{!G})$؛ او (ب)
$\Th{T}
\Proves !G$ له
$\Th{T} \Proves \lnot \OProv[\Th{T}](\gn{!G})
\liff !G$ سره يوځاے دا راکوي چې
$\Th{T} \Proves \lnot \OProv[\Th{T}](\gn{!G})$۔ په دې صورت کښې
به~$\Th{T}$ ناسازګاره وي۔ اوس طبيعي پوښتنه دا ده چې د
$\OProv[\Th{T}](y)$ د ثابت ټکي حالت څه دے؛ يعنې د داسې
!!a{sentence}~$!H$ حالت چې
\[
\Th{T} \Proves \OProv[\Th{T}](\gn{!H}) \liff !H.
\]
که دا جمله !!{derivable} واے، د شرط~(۱) له مخې به
$\Th{T} \Proves \OProv[\Th{T}](\gn{!H})$ و؛ خو همدا پايله د
پورته معادلۍ د مودس پوننس په کارولو هم ترلاسه کېږي۔ نو لږ تر لږه
د ګوډل د جملې د پخواني استدلال له لارې د~$\Th{T}$ ناسازګاري نۀ
راځي۔ دا، البته، نۀ ښيي چې~$\Th{T}$ په رښتيا~$!H$
!!{derive} \emph{کوي}۔

که پوښتنه لږ عمومي کړو، پرمختګ کولے شو۔ د ثابت ټکي د معادلۍ
له کيڼ څخه ښي ته لوری،
$\OProv[\Th{T}](\gn{!H}) \lif !H$، د يوې عمومي شېما بېلګه ده
چې د \emph{انعکاس اصل} بلل کېږي:
$\OProv[\Th{T}](\gn{!A}) \lif !A$۔ دا نوم ځکه ورکول کېږي چې
په يوه معنا ښيي~$\Th{T}$ د هغو څه په باب «انعکاس» کولے شي چې
!!{derive} کوي؛ په بنسټيز ډول د هر~$!A$ دپاره وايي: «که~$\Th{T}$،
$!A$ !!{derive} کولے شي، نو~$!A$ رښتيا دے»۔ دا خبره، البته،
يوازې د سمې تيورۍ دپاره رښتيا ده؛ نو داسې ښکاري چې تيورۍ به
په عمومي توګه د دې شېما ټولې بېلګې نۀ !!{derive} کوي۔ نو يوه
کافي پياوړې تيوري چې د !!{derivability} شرطونه پوره کوي، کومې
بېلګې !!{derive} کولے شي؟ هغه ټولې بېلګې چې~$!A$ پخپله
!!{derivable} وي، يقيناً۔ لاندې پايله ښيي چې بس همدومره دي۔

\begin{thm}\ollabel{thm:lob}
فرض کړئ~$\Th{T}$ داسې !!a{axiomatizable} تيوري ده چې~$\Th{Q}$
غځوي، او~$\OProv[\Th{T}](y)$ هغه فورمول دے چې د
\olref[2in]{sec} شرطونه P1--P3 پوره کوي۔ که~$\Th{T}$،
$\OProv[\Th{T}](\gn{!A}) \lif !A$ !!{derive}s کړي، نو په حقيقت
کښې~$\Th{T}$، $!A$ هم !!{derive}s کوي۔
\end{thm}

په بل ډول: که $\Th{T} \Proves/ !A$ وي، نو
$\Th{T} \Proves/
\OProv[\Th{T}](\gn{!A}) \lif !A$۔ دا پايله د لوب د قضيې په نوم
يادېږي۔

\begin{explain}
د لوب د قضيې د ثبوت لارښوده مفکوره د دې يو ځيرک «ثبوت» دے چې
سانتا کلاز شته۔ (که دغه پايله مو نۀ خوښېږي، هره بله خوښه پايله
پر ځاے کولے شئ۔) استدلال داسې دے:
\begin{enumerate}
\item $X$ دې دا جمله وي: «که~$X$ رښتيا وي، نو سانتا کلاز شته»۔
\item فرض کړئ~$X$ رښتيا ده۔
\item نو د هغې ويل شوې خبره برقرار ده؛ يعنې که~$X$ رښتيا وي،
  سانتا کلاز شته۔
\item ځکه چې فرض مو کړے~$X$ رښتيا ده، له (۲) او~(۳) څخه د
  مودس پوننس په کارولو پايله اخلو چې سانتا کلاز شته۔
\item موږ (۴)، «سانتا کلاز شته»، له (۲)، «$X$ رښتيا ده»، څخه
  وايستله۔ د شرطي ثبوت له مخې مو وښودله: «که~$X$ رښتيا وي، نو
  سانتا کلاز شته»۔
\item خو دا هماغه جمله~$X$ ده۔ نو ومو ښودله چې~$X$ رښتيا ده۔
\item بيا د پورته (۲)--(۴) استدلال له مخې سانتا کلاز شته۔
\end{enumerate}
که په دې مفکوره کښې «رښتيا ده» په «!!{derivable} ده» او «سانتا
کلاز شته» په~$!A$ واړوو او هغه صوري کړو، د لوب د قضيې ثبوت
ترې جوړېږي۔ چاره دا ده چې د ثابت ټکي لمه په
!!{formula}~$\OProv[\Th{T}](y) \lif !A$ پلې کړو۔ د دې فورمول
ثابت ټکی په پورته لنډ استدلال کښې له~$X$ جملې سره سمون لري۔
\end{explain}

\begin{proof}[د \olref{thm:lob} ثبوت]
فرض کړئ~$!A$ داسې !!a{sentence} ده چې~$\Th{T}$،
$\OProv[\Th{T}](\gn{!A}) \lif !A$ !!{derive}s کوي۔ $!B(y)$
دې !!{formula}~$\OProv[\Th{T}](y)
\lif !A$ وي، او د ثابت ټکي د لمې په کارولو داسې
!!a{sentence}~$!D$ ومومئ چې~$\Th{T}$،
$!D \liff !B(\gn{!D})$ !!{derive}s کوي۔ بيا لاندې هره يوه
په~$\Th{T}$ کښې !!{derivable} ده:
\begin{align}
  & !D \liff (\OProv[\Th{T}](\gn{!D}) \lif !A) \ollabel{L-1}\\
  & \qquad \text{$!D$ د~$!B(y)$ ثابت ټکی دے}\notag \\
  & !D \lif (\OProv[\Th{T}](\gn{!D}) \lif !A) \ollabel{L-2}\\
  & \qquad\text{له \olref{L-1} څخه}\notag\\
  & \OProv[\Th{T}](\gn{!D \lif (\OProv[\Th{T}](\gn{!D}) \lif !A)}) \ollabel{L-3}\\
  & \qquad \text{له \olref{L-2} څخه د P1 شرط له مخې}\notag \\
  & \OProv[\Th{T}](\gn{!D}) \lif \OProv[\Th{T}](\gn{\OProv[\Th{T}](\gn{!D}) \lif !A})
  \ollabel{L-4}\\
  &\qquad \text{له \olref{L-3} څخه د P2 شرط له مخې}\notag \\
  & \OProv[\Th{T}](\gn{!D}) \lif (\OProv[\Th{T}](\gn{\OProv[\Th{T}](\gn{!D})}) \lif \OProv[\Th{T}](\gn{!A})) \ollabel{L-5}\\
  &\qquad \text{له \olref{L-4} څخه د P2 شرط له مخې بيا} \notag\\
& \OProv[\Th{T}](\gn{!D}) \lif \OProv[\Th{T}](\gn{\OProv[\Th{T}](\gn{!D})}) \ollabel{L-6}\\
  & \qquad\text{د !!{derivability} د P3 شرط له مخې} \notag\\
  & \OProv[\Th{T}](\gn{!D}) \lif \OProv[\Th{T}](\gn{!A}) \ollabel{L-7} \\
  &\qquad\text{له \olref{L-5} او \olref{L-6} څخه}\notag\\
  & \OProv[\Th{T}](\gn{!A}) \lif !A \ollabel{L-8}\\
  &\qquad\text{د قضيې د فرض له مخې} \notag\\
  & \OProv[\Th{T}](\gn{!D}) \lif !A \ollabel{L-9}\\
  &\qquad\text{له \olref{L-7} او \olref{L-8} څخه}\notag\\
  & (\OProv[\Th{T}](\gn{!D}) \lif !A) \lif !D \ollabel{L-10}\\
  & \qquad \text{له \olref{L-1} څخه}\notag \\
  & !D \ollabel{L-11}\\
  & \qquad\text{له \olref{L-9} او \olref{L-10} څخه}\notag \\
  & \OProv[\Th{T}](\gn{!D}) \ollabel{L-12}\\
  & \qquad\text{له \olref{L-11} څخه د P1 شرط له مخې}\notag \\
  & !A \qquad\qquad\text{له \olref{L-8} او \olref{L-12} څخه}\notag
\end{align}
\end{proof}

د لوب قضيه چې ولرو، د ناتکميلۍ د دويمې قضيې لنډ ثبوت ترې راځي
(د هغو تيوريو دپاره چې د P1--P3 شرطونه پوره کوونکے
!!a{derivability} پريديکات لري): که
$\Th{T} \Proves
\OProv[\Th{T}](\gn{\lfalse}) \lif \lfalse$ وي، نو
$\Th{T} \Proves \lfalse$۔ که~$\Th{T}$ سازګاره وي،
$\Th{T} \Proves/ \lfalse$۔ نو
$\Th{T}
\Proves/ \OProv[\Th{T}](\gn{\lfalse}) \lif \lfalse$، يعنې
$\Th{T} \Proves/
\OCon[\Th{T}]$۔ له دې څخه دا هم ښودلے شو چې~$!H$، د
$\OProv[\Th{T}](x)$ ثابت ټکی، !!{derivable} دے۔ ځکه چې
\begin{align*}
  \Th{T} & \Proves \OProv[\Th{T}](\gn{!H}) \liff !H\\
  \intertext{په ځانګړي ډول}
    \Th{T} & \Proves \OProv[\Th{T}](\gn{!H}) \lif !H
\end{align*}
او ځکه د لوب د قضيې له مخې $\Th{T} \Proves !H$۔

% Going in the other direction, for homework I
% may ask you to work through a short proof of L\"ob's theorem, using
% the second incompleteness theorem instead of the fixed-point lemma.

\begin{prob}
فرض کړئ~$\Th{T}$ په محاسبوي ډول بديهي‌کړې تيوري ده، او
$\OProv[\Th{T}]$ د~$\Th{T}$ دپاره !!a{derivability} پريديکات
دے۔ لاندې څلور بيانونه وګورئ:
\begin{enumerate}
\item که $T \Proves !A$ وي، نو $T \Proves \OProv[\Th{T}](\gn{!A})$۔
\item $T \Proves !A \lif \OProv[\Th{T}](\gn{!A})$۔
\item که $T \Proves \OProv[\Th{T}](\gn{!A})$ وي، نو $T \Proves !A$۔
\item $T \Proves \OProv[\Th{T}](\gn{!A}) \lif !A$۔
\end{enumerate}
له کومو شرطونو سره له دې بيانونو هر يو رښتيا دے؟
\end{prob}

\end{document}
